\section{Case Study}
\label{app:case_study}

To qualitatively demonstrate the efficacy of DynaDebate, we analyze a representative failure case from the MATH500 dataset.

% =======================================================
% 1. 题目部分 (独立盒子)
% =======================================================
\begin{tcolorbox}[
    colback=white,
    colframe=black!70,
    boxrule=0.8pt,
    title=\textbf{Problem Example (MATH500)},
    fonttitle=\small\sffamily\bfseries,
    sharp corners,
    bottom=5pt
]
    \small
    \textbf{Question:} The expression $2\cdot 3 \cdot 4\cdot 5+1$ is equal to 121... In total, how many values can be obtained from the expression $2\cdot 3\cdot 4 \cdot 5 + 1$ by inserting parentheses? (Note that rearranging terms is not allowed, only inserting parentheses).

    \vspace{0.2cm}
    \textbf{Ground Truth:} 4 values ($144, 126, 122, 121$).
    \begin{itemize}[nosep, leftmargin=1em]
        \item $2\cdot 3 \cdot 4 \cdot (5 + 1) = 144$
        \item $2\cdot 3 \cdot (4 \cdot 5 + 1) = 126$
        \item $2\cdot (3 \cdot 4 \cdot 5 + 1) = 122$  \textcolor{red}{\textbf{(Often missed by baselines)}}
        \item $(2\cdot 3 \cdot 4 \cdot 5) + 1 = 121$
    \end{itemize}
\end{tcolorbox}

% =======================================================
% 2. 基线 Baseline (独立盒子 + breakable)
% =======================================================
\begin{tcolorbox}[
    breakable,
    colback=red!5!white,
    colframe=red!60!black,
    title=\textbf{\sffamily Baseline (SoM)},
    fonttitle=\small\bfseries,
    boxrule=0.5pt,
    sharp corners,
    left=2pt, right=2pt, top=2pt, bottom=2pt
]
    \footnotesize
    \textbf{Round 1:}
    \begin{itemize}[leftmargin=1.5em, nosep]
        \item \textbf{Agent A:} The expression is $2 \cdot 3 \cdot 4 \cdot 5 + 1$. By operator precedence, it evaluates to $120 + 1 = 121$.

Multiplication is associative, so any parentheses applied only to the product $2\cdot3\cdot4\cdot5$ still yield 120. Thus, all such groupings give 121.

However, if we group the addition early, e.g. $(5+1)$, new values arise:
\[
(2 \cdot 3) \cdot (4 \cdot (5 + 1)) = 6 \cdot 24 = 144.
\]
Enumerating all 14 valid parenthesizations shows only two distinct results: **121** and **144**.

\medskip
\noindent\textbf{Final Answer:} $\boxed{2}$

        \item \textbf{Agent B:} The base value is $(2 \cdot 3 \cdot 4 \cdot 5) + 1 = 121$.

Since multiplication is associative and commutative, re-grouping the product alone does not change the result.

Only when parentheses include the ``$+1$'' in a multiplicative subexpression do we get a new value:
\[
(2 \cdot (3 \cdot 4)) \cdot (5 + 1) = 24 \cdot 6 = 144.
\]
All other valid groupings either yield 121 or 144. No third value appears.

\medskip
\noindent\textbf{Final Answer:} $\boxed{2}$

        \item \textbf{Agent C:}We consider groupings where addition occurs at different depths:

\begin{align*}
(2 \cdot 3 \cdot 4 \cdot 5) + 1 &= 121, \\
(2 \cdot (3 \cdot 4)) \cdot (5 + 1) &= 144, \\
2 \cdot (3 \cdot ((4 \cdot 5) + 1)) &= 2 \cdot (3 \cdot 21) = 126.
\end{align*}

Other groupings replicate these values. Thus, three distinct results are found: 121, 144, and 126.

\medskip
\noindent\textbf{Final Answer:} $\boxed{3}$

    \end{itemize}

    \vspace{0.1cm}
    \textbf{Round 2 :}
    \begin{itemize}[leftmargin=1.5em, nosep]
        \item \textbf{Agent A :} The expression $2 \cdot 3 \cdot 4 \cdot 5 + 1$ evaluates to 121 by default.

There are 14 valid parenthesizations (4th Catalan number). Evaluating them yields:

\begin{align*}
\text{Most groupings} &\Rightarrow 121, \\
\text{Groupings with } (5+1) &\Rightarrow 144, \\
2 \cdot (3 \cdot ((4 \cdot 5) + 1)) &= 2 \cdot (3 \cdot 21) = 126.
\end{align*}

Thus, the distinct values are 121, 144, and 126.

\medskip
\noindent\textbf{Final Answer:} $\boxed{3}$

        \item \textbf{Agent B :} Although there are 42 parenthesizations (5th Catalan number), associativity of multiplication collapses most to the same value.

Key distinct evaluations:
\begin{align*}
(2 \cdot 3 \cdot 4 \cdot 5) + 1 &= 121, \\
(2 \cdot (3 \cdot 4)) \cdot (5 + 1) &= 144, \\
2 \cdot (3 \cdot ((4 \cdot 5) + 1)) &= 126.
\end{align*}

All other groupings reproduce one of these three values.

\medskip
\noindent\textbf{Final Answer:} $\boxed{3}$

        \item \textbf{Agent C:} Using recursive evaluation of all valid parenthesizations of $2 \cdot 3 \cdot 4 \cdot 5 + 1$, we find only three distinct results:

\[
121,\quad 144,\quad 126,
\]
corresponding to:
\begin{itemize}[left=0pt,topsep=2pt,itemsep=0pt]
    \item addition applied last,
    \item $5+1$ grouped early,
    \item $(4\cdot5 + 1)$ grouped before outer multiplications.
\end{itemize}

No other values arise under the no-rearrangement constraint.

\medskip
\noindent\textbf{Final Answer:} $\boxed{3}$
    \end{itemize}

    \vspace{0.1cm}
    \textbf{VOTED Answer:} $\boxed{3}$ \textcolor{red}{Incorrect}
\end{tcolorbox}

% =======================================================
% 3. 你的方法 DynaDebate (独立盒子 + breakable)
% =======================================================
\begin{tcolorbox}[
    breakable,
    lines before break=0,
    enlarge top initially by=0mm,
    pad at break=1mm,
    colback=blue!5!white,
    colframe=blue!60!black,
    title=\textbf{\sffamily DynaDebate (Ours)},
    fonttitle=\small\bfseries,
    boxrule=0.5pt,
    sharp corners,
    left=2pt, right=2pt, top=2pt, bottom=2pt
]
    \footnotesize
    \textbf{Phase 1: Path Allocation}
    \begin{itemize}[leftmargin=1.5em, nosep]
        \item \textbf{Path Generation Agent :}Three distinct solution strategies:

\begin{enumerate}[left=0pt,topsep=2pt,itemsep=0pt]
    \item \textbf{Combinatorics with Parsing Trees}: Model parenthesizations as binary trees; high reliability.
    \item \textbf{Dynamic Programming with Memoization}: Recursively compute all sub-expression values; high reliability.
    \item \textbf{Operator Precedence Analysis}: Manually inspect groupings; medium reliability for small expressions.
\end{enumerate}

    \end{itemize}

    \vspace{0.1cm}
    \textbf{Phase 2: Debate Round 1 }
    \begin{itemize}[leftmargin=1.5em, nosep]
        \item \textbf{Agent A:} Evaluating all 14 valid parenthesizations of $2 \cdot 3 \cdot 4 \cdot 5 + 1$ yields four distinct values:

\begin{align*}
(2\cdot3\cdot4\cdot5) + 1 &= 121, \\
2 \cdot (3 \cdot ((4 \cdot 5) + 1)) &= 126, \\
2 \cdot ((3 \cdot (4 \cdot 5)) + 1) &= 122, \\
\text{Any grouping with } (5+1) \text{ in product} &= 144.
\end{align*}

Thus, the number of distinct values is $\boxed{4}$.
        \item \textbf{Agent B:} Using dynamic programming on tokens $[2,\cdot,3,\cdot,4,\cdot,5,+,1]$, we compute:

\[
\texttt{dp}[0][4] = \{121,\ 144,\ 126\}.
\]

The split at the final '+' gives 121; splits before '+' with $(5+1)$ yield 144 or 126.
No 122 is generated due to state aggregation.

\medskip
\noindent\textbf{Final Answer:} $\boxed{3}$
        \item \textbf{Agent C:} Two cases only:
\begin{itemize}[left=0pt,topsep=2pt,itemsep=0pt]
    \item Addition last: $(2\cdot3\cdot4\cdot5)+1 = 121$,
    \item Addition grouped early (e.g., $5+1$): yields $144$.
\end{itemize}

All other groupings reproduce one of these. Thus, only two distinct values exist.

\medskip
\noindent\textbf{Final Answer:} $\boxed{2}$
    \end{itemize}

    \vspace{0.1cm}
    \textbf{Phase 3: Debate Round 2 }
    \begin{itemize}[leftmargin=1.5em, nosep]
        \item \textbf{Agent A:} Systematic evaluation of all valid parenthesizations of $2 \cdot 3 \cdot 4 \cdot 5 + 1$ yields four distinct values:

\begin{align*}
(2\cdot3\cdot4\cdot5) + 1 &= 121, \\
2 \cdot (3 \cdot ((4 \cdot 5) + 1)) &= 126, \\
2 \cdot ((3 \cdot (4 \cdot 5)) + 1) &= 122, \\
\text{Any grouping with } (5+1) \text{ in product} &= 144.
\end{align*}

Thus, the total number of distinct values is $\boxed{4}$.
        \item \textbf{Agent B :} Despite using dynamic programming, the evaluation only considered two outcomes:
\begin{itemize}[left=0pt,topsep=2pt,itemsep=0pt]
    \item Standard precedence: $121$,
    \item Grouping with $(5+1)$: $144$.
\end{itemize}
Other valid groupings (e.g., $2 \cdot (3 \cdot (20 + 1)) = 126$) were missed due to incomplete state enumeration.

\medskip
\noindent\textbf{Final Answer:} $\boxed{2}$
        \item \textbf{Agent C:} A complete dynamic programming evaluation over tokens $[2,\cdot,3,\cdot,4,\cdot,5,+,1]$ yields the full set of possible values:

\[
\{121,\ 144,\ 126,\ 122\}.
\]

These correspond to:
\begin{itemize}[left=0pt,topsep=2pt,itemsep=0pt]
    \item addition last,
    \item $(5+1)$ grouped early,
    \item $(4\cdot5 + 1)$ grouped,
    \item $(3\cdot4\cdot5 + 1)$ grouped.
\end{itemize}

All four are valid and distinct under the no-rearrangement constraint.

\medskip
\noindent\textbf{Final Answer:} $\boxed{4}$
    \end{itemize}

    \vspace{0.1cm}
    \textbf{VOTED Answer:} $\boxed{4}$ \textcolor{green}{Correct}

\end{tcolorbox}

\captionof{figure}{Comparison on a MATH500 problem involving operator precedence.In SoM, all agents overlooked the answer 122, whereas in DynaDebate, by assigning diverse reasoning paths, the agents obtained the correct answer (4); through structured debate, the influence of the correct answer was amplified and ultimately confirmed via voting.}
\label{fig:case_study_math}
